
\documentclass[12pt,a4paper]{article}
\usepackage[utf8]{inputenc}
\usepackage{graphicx}
\usepackage{color}
\usepackage[T1]{fontenc}
\usepackage[magyar]{babel}
\usepackage{listings}
\usepackage{float}

\usepackage{enumerate}
\usepackage{tikz}
\usetikzlibrary{shapes,arrows}
\usetikzlibrary{positioning}
\usetikzlibrary{arrows.meta}

\tikzset{
block/.style = {draw, fill=white, rectangle, minimum height=3em, minimum width=3em},
tmp/.style  = {coordinate}, 
input/.style = {coordinate},
output/.style= {coordinate},
pinstyle/.style = {pin edge={to-,thin,black}
}
}
\usepackage[siunitx]{circuitikz} 
\renewcommand{\lstlistingname}{Lista}
\usepackage{hyperref}
\hypersetup{
	pdftitle={HA5KFU tanfolyam - Szoftverrádió gyakorlat},
	pdfauthor={HA5KFU Rádióamatőr Klub},
	pdfsubject={HA5KFU tanfolyam},
	pdfcreator={latex},
	pdfkeywords={ },
	pdfpagemode=UseOutlines,
	pdfdisplaydoctitle=true,
	pdflang=hu,
	unicode
}
\usepackage{color} %red, green, blue, yellow, cyan, magenta, black, white
\usepackage{xcolor}
\usepackage{textcomp}

\definecolor{comment}{RGB}{0,128,0} % dark green
\definecolor{string}{RGB}{255,0,0}  % red
\definecolor{keyword}{RGB}{0,0,255} % blue

\usepackage{accsupp}
\newcommand{\noncopynumber}[1]{%
    \BeginAccSupp{method=escape,ActualText={}}%
    #1%
    \EndAccSupp{}%
}

\lstset{
	basicstyle=\footnotesize\ttfamily\color{black},
	commentstyle=\color{comment},
	stringstyle=\color{string},
	keywordstyle=\color{keyword},
	basicstyle=\footnotesize\ttfamily,
	numbers=left,
	numberstyle=\tiny\noncopynumber,
	numbersep=5pt,
	frame=lines,
	upquote=true,
	breaklines=true,
	prebreak=\raisebox{0ex}[0ex][0ex]{\ensuremath{\hookleftarrow}},
	showstringspaces=false,
	upquote=true,
	tabsize=2,
	columns=flexible,
}
\pagestyle{plain}
\sloppy
\begin{document}
\begin{center}
\includegraphics[width=300pt,keepaspectratio]{figures/ha5kfu.eps}
\\[0.5cm]
Rádióamatőr tanfolyamot segítő jegyzet, egyelőre kidolgozás alatt \\
Szabó Áron HA1FLX, HA5KFU... % Feel free to add yourself
\\
Retzler András HA7ILM szoftverrádiós vevője alapján
\\[1cm]

{\huge \bfseries Szoftverrádió gyakorlat megoldás \\[2cm]}



\end{center}
\subsection{Deriválás}

$\arctan$ deriválási szabálya:
\begin{equation}
\left( \arctan \left( x \right) \right)' = \frac{1}{1 + x^2}
\end{equation}
$\frac{f}{g}$ deriválási szabálya:
\begin{equation}
\left( \frac{f}{g} \right)' = \frac{ f' \cdot g - f  \cdot  g'}{g^2}
\end{equation}
$f(g(x))$ deriválási szabálya:
\begin{equation}
\left( f(g(x)) \right)' = f'(g(x)) \cdot g'(x)
\end{equation}
\vspace{20pt}

$\arctan \left( \frac{q}{i} \right)$ deriválása:
\begin{equation}
\frac{\partial}{\partial t}\left( \arctan \left( \frac{q}{i} \right) \right) = 
\end{equation}
\begin{equation}
\frac{1}{1 + \left( \frac{q}{i} \right)^2} \cdot \frac{\frac{\partial q}{\partial t} \cdot i - q  \cdot \frac{\partial i}{\partial t}}{i^2} =
\end{equation}
\begin{equation}
\frac{\frac{\partial q}{\partial t} \cdot i - q  \cdot \frac{\partial i}{\partial t}}{\left( 1 + \frac{q^2}{i^2} \right) \cdot i^2 } 
=
\end{equation}
\begin{equation}
\frac{i \cdot \frac{\partial q}{\partial t} - q \cdot \frac{\partial i}{\partial t}}{i^2 + q^2}
\end{equation}

A deriváltak közelítése diszkrét időben:

\begin{equation}
\frac{\partial x}{\partial t} \approx \frac{\Delta x}{\Delta t} = \frac{x[k]-x[k-1]}{T_s}
\end{equation}

\newpage
A kész kód így:

\begin{lstlisting}[frame=single,language=c,caption=FM demodulátor C-ben]
#include <math.h>
#include <stdio.h>


int main ()
{
    double i, q, s, ip, qp;
    // ip, qp az elozo i es q ertekek
    for (;;) // vegtelen ciklus
    {
        // beolvassuk az I mintat , az offszetet levonjuk , hogy a 0 tenyleg 0 legyen
        i =(( unsigned char ) getchar () -127) ;
        // Q- val ugyan ez
        q =(( unsigned char ) getchar () -127) ; // beolvassuk


        // di, dq a derivaltat kozeliti
        // diszkret idoben "igy kell derivalni"
        // csak egy konstans szorzo maradt le
        double di = i - ip;
        double dq = q - qp;

        // arctan derivaltja szerint, a szorzo az elejen kiserlezetes alapjan lett ennyi
        s = 40 * (di*q - dq*i)/(i*i+q*q);

        // elozo mintak beallitasa
        qp = q;
        ip = i;

        // s-t visszaalakitjuk offszetese majd kiirjuk
        putchar (( unsigned char ) ( s +127) ) ;
    }
}
\end{lstlisting}

\end{document}
